% §16 Exercises
%
% Cross-reference resolution status:
%
% DEFINED here:
%   \label{ex:ch16-1} through \label{ex:ch16-11}
%
% RESOLVED (used here, defined elsewhere):
%   \Cref{sec:varying-rates}                    -- §16.2
%   \Cref{sec:correlated-noise}                 -- §16.3
%   \Cref{sec:adversarial-inputs}               -- §16.4
%   \Cref{sec:where-framework-helps}            -- §16.5
%   \Cref{prop:azuma-hoeffding}                 -- §16.2
%   \Cref{prop:mcdiarmid}                       -- §16.3
%   \Cref{ax:element-indep}                     -- chapter 4 Axiom 1
%   \Cref{ax:block-rates-indep}                 -- chapter 4 Axiom 2
%   \Cref{thm:composition}                      -- chapter 3 composition theorem
%   \Cref{thm:chernoff}                         -- chapter 11 Chernoff bound
%   \Cref{def:bloom-filter}                     -- chapter 5 Bloom filter definition
%   \Cref{def:ppv}                              -- chapter 7 PPV definition
%   \eqref{eq:bloom-fpr}                        -- chapter 5 FPR formula
%
% 11 exercises, 3 star / 5 starstar / 3 starstarstar:
% - star x 3: identify-violation, Azuma-Hoeffding numerical, Bloom at capacity
% - starstar x 5: notebook-tied correlated noise, varying-rate composition,
%   Bloom adversary lower bound, Chernoff vs Azuma at iid, PPV under
%   varying prevalence
% - starstarstar x 3: McDiarmid for Bloom, adaptive adversary, sufficient
%   condition for composition theorem upper bound

% Difficulty-marker macros (matching chapters 1-15):
\providecommand{\exstar}{$\star$}
\providecommand{\exstarstar}{$\star\star$}
\providecommand{\exstarstarstar}{$\star\star\star$}

\section*{Exercises}
\addcontentsline{toc}{section}{Exercises}

\begin{enumerate}

% -----------------------------------------------------------------------
% WARM-UP  (star x 3): exercises 1, 2, 3
% -----------------------------------------------------------------------

\item \exstar
\label{ex:ch16-1}
\textbf{Identify the violation.}
For each of the following three systems, state which axiom is
violated (\Cref{ax:element-indep}, Axiom~1; \Cref{ax:block-rates-indep},
Axiom~2; or both) and which of \Cref{sec:varying-rates},
\Cref{sec:correlated-noise}, or \Cref{sec:adversarial-inputs}
addresses the failure.
\begin{enumerate}
  \item[(a)] A classifier whose training distribution drifts from
    the deployment distribution: the per-query error rate $\fprate(t)$
    is no longer stable across queries.
  \item[(b)] Two Bloom filters with a shared hash family, queried
    in parallel on the union of their designated sets.
  \item[(c)] An online-learning classifier that updates its weights
    after each query, so the noise at query $t$ depends on the
    responses at queries $1, \ldots, t - 1$.
\end{enumerate}
One-line answer per part is sufficient.

\item \exstar
\label{ex:ch16-2}
\textbf{Azuma-Hoeffding numerical application.}
For a martingale $(X_t)_{t = 0, 1, \ldots, n}$ with bounded
differences $|X_t - X_{t-1}| \le 1$ and $n = 100$, apply
\Cref{prop:azuma-hoeffding} to compute the upper bound on
$\Pr[|X_n - \mathbb{E}[X_n]| \ge 10]$. Plug $a = 10$ and
$\sum_{t = 1}^{n} c_t^2 = n = 100$ into the bound $2 \exp(-a^2 / (2
\sum c_t^2))$ and report both the exponent and the resulting
probability bound. State in one sentence what the bound means
operationally: roughly, how often the martingale should be expected
to wander at least $10$ units from its mean over the $100$-step run.

\item \exstar
\label{ex:ch16-3}
\textbf{Bloom filter at capacity.}
A Bloom filter is sized for $n = 1000$ items at target $\fprate =
0.01$, using the \eqref{eq:bloom-fpr} formula to fix the bit-array
size $m$ and the number of hash functions $k$. The filter is then
filled with $n_{\text{actual}} = 5000$ items rather than the design
$n = 1000$.
\begin{enumerate}
  \item[(a)] From the \eqref{eq:bloom-fpr} formula $\fprate \approx
    (1 - e^{-k n / m})^k$, recover the design $m$ and $k$ at $n =
    1000$ and $\fprate = 0.01$ (the chapter~5 standard
    construction gives $m / n \approx 9.6$ and $k \approx 7$).
  \item[(b)] Recompute $\fprate$ with the same $m$ and $k$ but
    $n = 5000$. The over-filling pushes the bit-array density past
    its design value; report the new $\fprate$ to two significant
    figures.
  \item[(c)] State which axiom is violated and which section
    diagnoses the failure.
\end{enumerate}

% -----------------------------------------------------------------------
% MEDIUM  (starstar x 5): exercises 4, 5, 6, 7, 8
% -----------------------------------------------------------------------

\item \exstarstar
\label{ex:ch16-4}
\textbf{Notebook-tied correlated-noise simulation.}
Use the chapter~16 Python notebook to simulate two Bloom filters
with a shared hash family.
\begin{enumerate}
  \item[(a)] Construct two designated sets $A$ and $B$ of $1000$
    random 32-bit integers each (disjoint), and build two Bloom
    filters of identical parameters $(m, k)$ at target $\fprate =
    0.01$, sharing the same hash family across both filters.
  \item[(b)] Generate $10^5$ query points drawn from outside $A
    \cup B$, and measure the joint event rate $\Pr[\text{filter}_A
    \text{ says } \top \text{ and filter}_B \text{ says } \top]$.
  \item[(c)] Compare the empirical joint rate to the iid prediction
    $\fprate_A \cdot \fprate_B \approx 10^{-4}$. Report the gap and
    state in one sentence which direction the shared-hash-family
    correlation pushes the joint rate (above the product, below,
    or unchanged) and why.
\end{enumerate}

\item \exstarstar
\label{ex:ch16-5}
\textbf{Composition theorem under varying rates.}
The chapter~3 composition theorem \Cref{thm:composition} gives
$\eta_{\text{total}} = 1 - \prod_{i = 1}^{k} (1 - \eta_i)$ for $k$
independent stages with fixed rates $\eta_i$. Suppose now that the
$\eta_i$ are themselves independent random variables with means
$\mu_i = \mathbb{E}[\eta_i]$ and variances $\sigma_i^2$.
\begin{enumerate}
  \item[(a)] Show that $\mathbb{E}[\eta_{\text{total}}] \le 1 -
    \prod_{i = 1}^{k} (1 - \mu_i)$, with strict inequality unless
    every $\sigma_i^2 = 0$. (Hint: $\mathbb{E}[1 - \eta_i] =
    1 - \mu_i$, and Jensen's inequality applied to the convex
    function $f(x) = -\log(1 - x)$ on the product.)
  \item[(b)] Quantify the gap for the symmetric case $\mu_i = \mu$
    and $\sigma_i^2 = \sigma^2$ for all $i$, with $k = 5$ stages,
    $\mu = 0.1$, and $\sigma = 0.02$. Report the deterministic
    bound, the expected combined rate, and the gap in percentage
    points.
  \item[(c)] State which axiom is violated and which section
    addresses the failure.
\end{enumerate}

\item \exstarstar
\label{ex:ch16-6}
\textbf{Bloom adversary lower bound.}
A Bloom filter is configured at $\fprate = 0.001$ under iid
queries from an oblivious source. An attacker is given the hash
family used in the filter's construction (a non-cryptographic
salt is published, for example).
\begin{enumerate}
  \item[(a)] Derive the expected number of queries the attacker
    needs to find an input $x \notin A$ for which the filter
    returns $\top$. (Hint: the attacker can hash candidates
    locally and predict which would collide; in the worst case
    they need to try $O(1 / \fprate)$ candidates before finding
    one that hits the colliding-bit-pattern set.)
  \item[(b)] Compare to the iid baseline expectation of $1 /
    \fprate = 1000$ random queries to find a false positive.
    Discuss in one paragraph how the attacker's gain depends on
    knowledge of the hash family.
  \item[(c)] State the defense: if the hash family is keyed by a
    secret salt unknown to the attacker, the per-query attacker
    cost reverts to $1 / \fprate$. Identify which section's
    catalogue of defenses this falls under.
\end{enumerate}

\item \exstarstar
\label{ex:ch16-7}
\textbf{Chernoff vs Azuma at iid.}
The chapter~11 multiplicative Chernoff bound \Cref{thm:chernoff}
applies to sums of iid Bernoulli random variables and gives an
exponential tail bound. The Azuma-Hoeffding inequality
\Cref{prop:azuma-hoeffding} applies to bounded-difference
martingales and is more general.
\begin{enumerate}
  \item[(a)] For a sum $S_n = \sum_{t = 1}^{n} X_t$ of iid Bernoulli
    random variables $X_t \in \{0, 1\}$ with mean $\mu = \mathbb{E}[X_t]$,
    verify that the partial-sum sequence $M_t = S_t - t \mu$ is a
    martingale with respect to its natural filtration, with bounded
    differences $|M_t - M_{t-1}| \le 1$.
  \item[(b)] Apply \Cref{prop:azuma-hoeffding} to $M_n$ with $a =
    \epsilon n$ for a deviation parameter $\epsilon > 0$. Show that
    the resulting bound $2 \exp(-\epsilon^2 n / 2)$ is of the same
    exponential order in $n$ as the Chernoff bound (up to
    multiplicative constants in the exponent), recovering the iid
    tail behavior.
  \item[(c)] State in one sentence the price the generality of
    Azuma-Hoeffding pays in the iid case: the constants in the
    exponent are weaker than the tighter Chernoff bound, even
    though the order in $n$ matches.
\end{enumerate}

\item \exstarstar
\label{ex:ch16-8}
\textbf{PPV under varying prevalence.}
The chapter~7 \Cref{def:ppv} expressed PPV as a function of
$\fprate$, $\fnrate$, and prevalence $\pi$. Suppose now that
$\pi(t)$ varies with the query index while $\fprate$ and
$\fnrate$ remain stable.
\begin{enumerate}
  \item[(a)] Write PPV at query $t$ as $\mathrm{PPV}(t) = \pi(t)
    (1 - \fnrate) / [\pi(t)(1 - \fnrate) + (1 - \pi(t)) \fprate]$,
    the formula of \Cref{def:ppv} evaluated at the per-query
    prevalence.
  \item[(b)] Derive the time-average PPV $\bar{\mathrm{PPV}} =
    (1 / m) \sum_{t = 1}^{m} \mathrm{PPV}(t)$ and the PPV
    evaluated at the time-average prevalence $\mathrm{PPV}(\bar\pi)$,
    where $\bar\pi = (1 / m) \sum \pi(t)$. Show by a small numerical
    example with $\pi(t) \in \{0.1, 0.5\}$ at equal frequencies,
    $\fprate = \fnrate = 0.05$, that $\bar{\mathrm{PPV}} \ne
    \mathrm{PPV}(\bar\pi)$ in general, and state which is larger.
  \item[(c)] Identify a sufficient condition under which the two
    quantities coincide. (Hint: PPV is linear in $\pi$ only in
    degenerate cases; otherwise Jensen's inequality applies.)
    Reference \Cref{sec:where-framework-helps}'s rule that PPV
    survives per-query even when iid fails.
\end{enumerate}

% -----------------------------------------------------------------------
% CHALLENGING / EXTENSION  (starstarstar x 3): exercises 9, 10, 11
% -----------------------------------------------------------------------

\item \exstarstarstar
\label{ex:ch16-9}
\textbf{McDiarmid for Bloom filter empirical FPR.}
Consider a Bloom filter from \Cref{def:bloom-filter} fixed at
build time, queried against a fixed (non-iid) stream of $n$ query
points $x_1, \ldots, x_n$. Let $\hat{\fprate} = (1 / n) \sum_{t = 1}^{n}
\mathbb{1}[\text{filter says } \top \text{ at } x_t]$ be the empirical
FPR.
\begin{enumerate}
  \item[(a)] Apply \Cref{prop:mcdiarmid} with the function $f(h_1,
    \ldots, h_k) = \hat{\fprate}$ viewed as a function of the
    random hash-family choices $h_1, \ldots, h_k$. Bound the
    per-coordinate influence: changing one hash function affects
    at most a fraction $1 / k$ of bits the filter sets, so the
    bounded-difference constant for each coordinate is
    $c_i \le 1$ (every empirical-FPR value is in $[0, 1]$, the
    crudest bound).
  \item[(b)] Apply \Cref{prop:mcdiarmid} to get a concentration
    bound for $\hat{\fprate}$ around its expectation
    $\mathbb{E}[\hat{\fprate}]$: state the resulting inequality
    of the form $\Pr[|\hat{\fprate} - \mathbb{E}[\hat{\fprate}]|
    \ge \epsilon] \le 2 \exp(-2 \epsilon^2 / (\sum c_i^2))$.
  \item[(c)] Compare to the iid Chernoff bound that would apply
    if the query stream were iid: identify the cost of the
    non-iid analysis in the bound's tightness, and discuss what
    a tighter influence bound (one that exploits the structure of
    the hash family) would buy.
\end{enumerate}

\item \exstarstarstar
\label{ex:ch16-10}
\textbf{Adaptive adversary.}
A non-adaptive attacker submits a fixed batch of $q$ queries to
a Bloom filter at $\fprate = \epsilon$ and observes the responses.
An adaptive attacker submits queries one at a time, using the
response to each query to choose the next.
\begin{enumerate}
  \item[(a)] State the non-adaptive baseline: under iid queries
    the expected number of false positives in $q$ queries is
    $q \epsilon$, and the expected number of queries to find the
    first false positive is $1 / \epsilon$.
  \item[(b)] Design an adaptive query strategy that exploits
    response feedback. (Suggestion: enumerate candidate inputs
    in an order parameterized by their hash-signature distance
    from previously observed positives; cluster candidates around
    known-colliding bit patterns to amplify the per-query hit
    rate.)
  \item[(c)] Quantify the adaptive gain over the non-adaptive
    baseline: under the hash-distance enumeration of (b), the
    adaptive attacker's expected queries-to-first-positive can
    drop below $1 / \epsilon$ by a factor depending on the
    correlation structure of the hash family. Derive a bound on
    the gain for $k$ independent hash functions and discuss the
    limit $k \to \infty$.
\end{enumerate}

\item \exstarstarstar
\label{ex:ch16-11}
\textbf{Bernoulli framework's domain of validity: composition
upper bound.}
The chapter~3 composition theorem \Cref{thm:composition}
$\eta_{\text{total}} = 1 - \prod_{i = 1}^{k} (1 - \eta_i)$ is an
equality under \Cref{ax:element-indep} (Axiom~1, independence)
and \Cref{ax:block-rates-indep} (Axiom~2, stable per-block rates).
State a sufficient condition under which the formula gives an
\emph{upper bound} on the true combined error rate, rather than
an equality.
\begin{enumerate}
  \item[(a)] Suppose the per-stage errors $\eta_i$ are not
    independent but are \emph{negatively correlated} or
    \emph{negatively associated}: for every pair $(i, j)$,
    $\Pr[\text{error at } i \text{ and at } j] \le \Pr[\text{error
    at } i] \cdot \Pr[\text{error at } j]$. Show that the
    composition formula then gives an upper bound rather than an
    equality.
  \item[(b)] Give one concrete example of a system that satisfies
    the negative-association condition. (Hint: a routing system
    that explicitly avoids re-using a failed component, so
    consecutive stages reuse less; or a sampling scheme that
    enforces a hard quota on errors per batch.)
  \item[(c)] Contrast with the positively correlated case (the
    shared-hash-family Bloom-filter scenario of
    \Cref{sec:correlated-noise}): show that under positive
    correlation, the composition formula gives a \emph{lower
    bound} on the true combined error rate, so the chapter~3
    inequality changes direction depending on the sign of the
    correlation. State the practical implication for a designer
    who knows only the marginal rates: the iid formula sits
    between two error-rate cases that bracket the true value.
\end{enumerate}

\end{enumerate}
